\documentclass[10pt,a4paper]{amsart}
\usepackage[margin=19mm]{geometry}
\usepackage{amsmath,amssymb,mathtools}
\usepackage[scr=rsfs,cal=euler]{mathalfa}
\usepackage{xfrac}
\usepackage[unicode,hidelinks,bookmarks=false]{hyperref}
\newcommand{\ie}{i.e.\ }
\newcommand{\cf}{cf.\ }
\newcommand{\resp}{respectively\ }
\newcommand{\Subsection}{\subsection}
\providecommand{\nobreakdash}{\nobreak\hbox{-}}
\setlength{\emergencystretch}{2em}
\multlinegap0pt
\usepackage{tikz}
\usetikzlibrary{cd}
\usepackage{tikz}
\usepackage{xcolor}
\usepackage{csquotes}
\usepackage{amssymb}
\newtheorem{proposition}{Proposition}
\newcommand{\DD}{\mathcal{D}}
\newcommand{\PP}{\mathcal{P}}
\newcommand{\ZZ}{\mathcal{Z}}
\DeclareMathOperator*{\colim}{colim}
\DeclareMathOperator*{\hocolim}{hocolim}
\title{The homotopy type of the moment-angle complex associated to the complex of injective words}
\author{Pedro Conceição}
\date{Source: arXiv:2603.09491v3 (CC BY 4.0)}
\begin{document}
\maketitle

\noindent\emph{Source and licence.} The homotopy type of the moment-angle complex associated to the complex of injective words. Version arXiv:2603.09491v3 (CC BY 4.0), distributed by arXiv under the Creative Commons Attribution 4.0 International licence, http://creativecommons.org/licenses/by/4.0/, as recorded in the arXiv licence metadata for this exact version and shown on the abstract page https://arxiv.org/abs/2603.09491v3. Full licence notice: \texttt{licenses/arxiv-2603.09491-CC-BY-4.0.txt}.\par\smallskip

\noindent\textit{Editorial scope.} Counted unit: the article's proposition push-glue (source0.tex lines 652--654). The article's complete original proof of it is delivered with it (source0.tex lines 655--672, one \texttt{proof} environment of 18 lines). No mathematics is written, repaired or re-derived: the statement and the proof are the article's own lines, in the article's own notation, and no retained line is substituted or re-flowed.  Retained uncounted as the source-local context the proof needs: lines 424--426; lines 584--600; lines 644--646.  Nothing was omitted from any retained span, so the package carries no omission record.  No retained line contains a citation, so the wrapper carries no bibliography and no bibliography entry.  The retained lines contain the article's own diagram code, which the wrapper typesets with the standard tikz packages; no image file is delivered and the package declares no asset dependency.  Editorial formatting only, in addition: an AMS \texttt{amsart} preamble that restates the article's own declarations and macros used by the retained lines, namely the environments \texttt{proposition} on separate counters, together with the standard packages those lines need.  The retained environments print in this document as Proposition 1 in order of appearance, whereas the article prints the same items with the numbering of its own full document.  The complete native source of the article as distributed in the arXiv e-print for this exact version is bundled unchanged as \texttt{sources/196066/source0.tex}.
\begin{equation}\label{union-equiv-levi}
	\ZZ_\PP\textbf{(X,A)} \simeq \hocolim_{I_\ZZ} \DD(\ZZ_\sigma) \simeq \colim_{I_\ZZ} \DD(\ZZ_\sigma) \simeq \bigcup_{\sigma \in \PP} Z_\sigma,
\end{equation}

\begin{proposition}\label{induce-pushout}
	Let $K$ be a finite ordered simplicial complex on the vertex set $[m]$. Suppose that there is a pushout of ordered simplicial complexes 
	\[
	\begin{tikzcd}
		I \arrow[r] \arrow[d] & K_2 \arrow[d] \\
		K_1 \arrow[r]  & K. 
	\end{tikzcd}
	\]
	Denote by $K, \bar I, \bar K_1, \bar K_2$ the associated face posets on the same vertex set of $K$. Let $\textbf{(X,A)} = \{(X_i,A_i)\}_{i=1}^m$ be a collection of pairs of pointed $CW$-complexes and consider the polyhedral products spaces $\ZZ_{\bar I}\textbf{(X,A)}$, $\ZZ_{\bar K_2}\textbf{(X,A)}$, $\ZZ_{\bar K_1}\textbf{(X,A)}$ and $\ZZ_{K}\textbf{(X,A)}$. Then, the diagram below
	\[
	\begin{tikzcd}
		\ZZ_{\bar I}\textbf{(X,A)} \arrow[r] \arrow[d] & \ZZ_{\bar K_2}\textbf{(X,A)} \arrow[d] \\
		\ZZ_{\bar K_1}\textbf{(X,A)} \arrow[r]  & \ZZ_{K}\textbf{(X,A)},
	\end{tikzcd}
	\] 
	is a homotopy pushout of polyhedral products. By the homotopy equivalence in \ref{union-equiv-levi}, namely $\ZZ_{K}\textbf{(X,A)} \simeq \bigcup_{\sigma \in K} Z_\sigma$, we conclude that $\ZZ_{K}\textbf{(X,A)} \simeq \ZZ_{{\bar K_1}}\textbf{(X,A)}\cup_{\ZZ_{\bar I}\textbf{(X,A)}} \ZZ_{\bar K_2}\textbf{(X,A)}$.  
\end{proposition}

\begin{proposition}\label{prop-pute-boundary-momentangle}
Let $I$ be a pure (i.e., all facets have same dimension) simplicial complex which is a codimension $1$ subcomplex of the boundary of the $(n-1)$-dimensional simplex $\Delta^{n-1}$. Assume further that $I$ contains all the $n$ vertices of $\Delta^{n-1}$. Then, $\ZZ_I \simeq S^{2|I|-1}$, where $|I|$ is the number of facets in $I$.
\end{proposition}

\begin{proposition}\label{push-glue}
	Let $K_1,K_2$ be two $(d-1)$-dimensional simplices $\Delta^{d-1}$ on same vertex set $[d]$ and $I$ be a pure codimension $1$ subcomplex formed by a non-empty set of facets of the boundary of $K_1$ (and $K_2$). Assume further that $I$ contains all the vertices $[d]$ of $\Delta^{d-1}$. Let $K = K_1 \cup_I K_2$ be the ordered simplicial complex obtained as pushout of $K_1$ attached to $K_2$ along $i: I \to K_2$ (i.e. attaching $I$ identically to the corresponding facets of $\partial K_2$). Then, $\ZZ_K \simeq S^{2|I|}$, where $|I|$ is the number of facets of $I$. 
\end{proposition}

\begin{proof}
	By hypothesis $K_1,K_2$ and $I$ have no ghost vertices. Therefore, Proposition \ref{induce-pushout} yields the following homotopy pushout of polyhedral products
	\[
	\begin{tikzcd}
		\ZZ_{I}\arrow[r, "\ZZ(i)"] \arrow[d] & \ZZ_{K_2} \simeq * \arrow[d] \\
		\ZZ_{K_1} \simeq * \arrow[r]  & \ZZ_{K},
	\end{tikzcd}
	\]
	where $\ZZ_{K_1} = \ZZ_{K_1} = (D^{2d}) \simeq *$ since both $K_1$ and $K_2$ are copies of a $(d-1)$-simplex $\Delta^{d-1}$. Moreover, Proposition \ref{prop-pute-boundary-momentangle} showed that $\ZZ_{I} \simeq S^{2|I|-1}$, where $|I|$ is the number of facets of $I$.
 	Hence, the homotopy pushout
	\[
	\begin{tikzcd}
		S^{2|I| -1}\arrow[r, "\ZZ(i)"] \arrow[d] &  (D^{2d}) \simeq * \arrow[d] \\
		(D^{2d}) \simeq *  \arrow[r]  & \ZZ_{K},
	\end{tikzcd}
	\]
     implying that $\ZZ_{K} \simeq \Sigma S^{2|I| -1} \simeq S^{2|I|}$.	
\end{proof}

\end{document}
